%% OREGON CIRCUIT COURT PLEADING PAPER (UTCR 2.100)
%% XeLaTeX required
\documentclass[12pt,letterpaper]{article}
\usepackage{xltxtra,fontspec}
\usepackage{geometry}
\usepackage{fancyhdr}
\usepackage{setspace}
\usepackage{enumitem}
\usepackage{titlesec}
\usepackage{hyperref}
\usepackage{amssymb}
\usepackage{array}
\usepackage{lineno}

\setmainfont{TeX Gyre Termes}

\geometry{
  letterpaper,
  left=1.75in,
  right=1in,
  top=1in,
  bottom=1in
}

\doublespacing
\setlength{\parindent}{0pt}
\setlength{\parskip}{6pt}

\pagestyle{fancy}
\fancyhf{}
\renewcommand{\headrulewidth}{0pt}
\fancyfoot[C]{\thepage}

\renewcommand{\linenumberfont}{\normalfont\small}
\setlength{\linenumbersep}{1.0in}
\linenumbers

\titleformat{\section}{\normalsize\bfseries\centering}{}{0em}{}
\titleformat{\subsection}{\normalsize\bfseries}{}{0em}{}
\titlespacing*{\section}{0pt}{12pt}{6pt}
\titlespacing*{\subsection}{0pt}{12pt}{6pt}

\hypersetup{colorlinks=false}

\begin{document}

%% ============================================================
%% CAPTION
%% ============================================================

\begin{center}
{\large\bfseries IN THE CIRCUIT COURT OF THE STATE OF OREGON}\\[4pt]
{\large\bfseries FOR THE COUNTY OF WASHINGTON}\\[18pt]
\end{center}

\noindent
\begin{tabular}{p{3.0in} p{0.1in} p{2.7in}}
STATE OF OREGON, & & \\[2pt]
\hspace{0.4in} Plaintiff, & & Case No.: \textbf{24CR49043}\\[6pt]
\hspace{0.2in} v. & & \textbf{NOTICE OF APPEARANCE PRO SE;}\\[2pt]
ASHLE\textquotesingle{} ANTOINETTE PENN, & & \textbf{NOTICE OF PRESERVATION OF}\\[2pt]
\hspace{0.4in} Defendant. & & \textbf{RECORD; NOTICE OF SYSTEMATIC}\\[2pt]
 & & \textbf{OBSTRUCTION OF ACCESS TO}\\[2pt]
 & & \textbf{PROCEEDINGS; AND OBJECTION}\\[2pt]
 & & \textbf{TO CHARACTERIZATION OF}\\[2pt]
 & & \textbf{DEFENDANT\textquotesingle S STATUS}\\[6pt]
 & & \textit{Special Appearance Only;}\\[2pt]
 & & \textit{Filed on the Papers}\\
\end{tabular}

\vspace{18pt}

%% ============================================================
%% SECTION I
%% ============================================================

\section*{I. NATURE AND PURPOSE OF THIS NOTICE}

Defendant Ashle\textquotesingle{} Antoinette Penn appears specially and pro se to
establish a complete and contemporaneous record on matters that have advanced, in
Defendant's absence and without her knowledge, in a manner inconsistent with her
constitutional and statutory rights. This Notice neither litigates nor waives any
substantive claim. It states what the record shows, names what the record requires
this Court to see, and places Defendant squarely before this Court on her own terms.

The record of these proceedings, taken as a whole, reflects a single coherent
pattern: Defendant has not absented herself from this Court. She has been
systematically denied access to it. That distinction is not rhetorical. It is the
organizing fact of this filing, and it will be the organizing fact of the motions
that follow.

This filing is a special appearance. Defendant does not submit to general
jurisdiction and does not waive any pending motion, jurisdictional challenge, or
constitutional objection by this filing.

%% ============================================================
%% SECTION II
%% ============================================================

\section*{II. DEFENDANT HAS BEEN PRESENT --- THE RECORD CONFIRMS IT}

Defendant appeared at every proceeding at which she had counsel. She appeared at
arraignment, at hearings before the Honorable Brandon Thompson on both case
24CR49043 and case 24CR27775, and at the April 28, 2025 hearing at which, over
her express objection stated on the record, her second appointed attorney was
removed from the case.

At the April 28, 2025 hearing, Defendant had independently prepared and filed both
a demurrer and a declaration because appointed counsel had represented she would
file a demurrer and did not. Defendant appeared. Defendant spoke. Defendant asked
to preserve the record. The Court told her she would not be prejudiced by the
deferral of argument on those filings that day. Defendant relied on that
representation.

At timestamp 19:05 through 19:09 of that same hearing, the court clerk confirmed
Defendant's current address on the record. Defendant followed the OPDC check-in
process as directed. She followed every instruction this Court gave her.

She has attempted to participate in subsequent proceedings by every means made
available to her. She has been blocked at each attempt by mechanisms this Notice
describes precisely.

\subsection*{A. Defendant Is Not Evading These Proceedings}

Deputy Rambin's December 29, 2025 affidavit in support of the arrest warrant in
case 25CN05790 states that multiple deputies have attempted to contact Defendant
at her last known address and have been unsuccessful, and that it ``appears that
she is actively attempting to avoid contact with law enforcement.''

The docket does not support that characterization. Defendant appeared when she
had counsel. When she had no counsel, she was told she would not have a court date
until new counsel was appointed. She received no adequate notice of subsequent
proceedings, for reasons this Notice describes in detail. She attempted to appear
remotely when she became aware of a hearing and was denied.

Defendant further notes that Deputy Rambin's ``last known address'' in his
December 29, 2025 affidavit is 17548 NW Springville Road, Apt. 9, Portland ---
the address at which law enforcement executed the May 10, 2025 warrant and at
which Defendant was located. That is not the address of a person avoiding law
enforcement. That is a person's home. The record does not support the
characterization of evasion, and Defendant objects to it.

%% ============================================================
%% SECTION III
%% ============================================================

\section*{III. DEFENDANT DID NOT RECEIVE NOTICE --- THE MECHANISM IS
DOCUMENTED}

\subsection*{A. OPDC Is an Administrative Placeholder, Not Representation}

At the April 28, 2025 hearing, the Honorable Brandon Thompson explained at
timestamps 17:28 through 17:50 that Defendant's case would move to what he
described as ``OPDC doc'' status --- a queue designation reflecting that no
attorney had been assigned. The Court stated explicitly: ``we don't have a lawyer
for you right now.'' Defendant was directed to call a check-in number and was
told she would not have another court date until a new lawyer was obtained.

The OPDC designation is an administrative queue. It is not an attorney-client
relationship. It does not constitute representation. It carries no obligation on
OPDC's part to monitor the docket, attend hearings, or transmit notices to
Defendant. Nevertheless, the court's automated notice system generated and routed
hearing notices to OPDC in Salem as though Defendant were represented. Those
notices did not reach Defendant.

The practical consequence requires no inference: proceedings were scheduled, held,
and concluded on matters directly affecting Defendant's liberty and the charges
against her, and Defendant received no notice of any of them.

\subsection*{B. The Address of Record Was Altered When the Contempt Case Was
Filed --- and That Alteration Propagated Across All Active Cases}

On December 30, 2025, case 25CN05790 --- the contempt action --- was filed. That
filing introduced a stale address for Defendant into the court's case management
system: an address at which Defendant has not resided for approximately ten years.
Because case 25CN05790 was the most recently filed case involving Defendant, the
court's system pulled that address forward and applied it as Defendant's address
of record in her other active cases, including 24CR49043 and 24CR27775. Defendant
was informed of this propagation mechanism by a court release officer.

The consequence was immediate and total: every notice generated in every active
case after December 30, 2025 was directed to an address at which Defendant does
not live and has not lived for approximately a decade.

This was not a routine clerical error. The address entered into case 25CN05790 on
December 30, 2025 was stale. It was entered into the most recently filed case.
The system did the rest. The timing is not incidental --- the contempt case was
filed, the stale address was entered, and the notice blackout began,
simultaneously, on the same date.

The internal contradiction that establishes the State's actual knowledge of
Defendant's current address is this: the arrest warrant in case 25CN05790 was
executed at 17548 NW Springville Road, Apt. 9, Portland --- Defendant's actual
current residence. The State executed process at Defendant's correct address
while the court system simultaneously routed notices to a different, incorrect
address. The State cannot simultaneously claim ignorance of Defendant's location
and have executed a warrant there.

\subsection*{C. A Court Release Officer Sent Notice to the Stale Address and
Asserted That Sending Constituted Delivery}

Christina Stiefer, Pretrial/Verification/Set Aside/Traffic Supervisor for
Washington County Circuit Court, transmitted a notice to Defendant at the stale
address of record. When Defendant raised the address problem, the officer
maintained that the notice had been properly sent --- treating the act of
transmission to a wrong address as equivalent to delivery. Defendant was
instructed to surrender herself notwithstanding the defective notice.

Proof of sending is not proof of receipt. An address that appears in the court
system as a result of a filing-date propagation mechanism does not become
Defendant's address by virtue of its appearance there. Defendant preserves all
rights arising from this defective notice.

\subsection*{D. Defendant's Motion to Appear Remotely Was Denied; Her Absence
Was Then Recorded as Failure to Appear}

When Defendant learned that proceedings were being held in case 24CR27775 in
connection with the court's sua sponte consideration of mandatory dismissal under
\textit{State v.\ Roberts}, 374 Or.\ 821 (2026), she filed a motion to appear
remotely pursuant to Presiding Judge Order No.\ 343 (Remote Proceedings), signed
February 15, 2024, by the Honorable Rebecca D.\ Guptill, Presiding Judge.

That motion was denied. Defendant was then noted as having failed to appear at
the proceeding from which she had been excluded by the denial of the only
mechanism by which she sought to attend.

The sequence is this: Defendant sought access. Access was denied. The denial was
then converted into a procedural finding against Defendant. This Court is entitled
to know that the failure-to-appear notation followed directly from a denial of
the means to appear.

%% ============================================================
%% SECTION IV
%% ============================================================

\section*{IV. ADVERSE PROCEEDINGS WERE CONDUCTED AFTER THE COURT
REPRESENTED THAT PROCEEDINGS WERE SUSPENDED}

At the April 28, 2025 hearing, at timestamp 18:18, the Honorable Brandon Thompson
stated on the record: \textit{``it cancels all future hearing dates at this
time.''} The Court further represented that Defendant would not have another court
date until new counsel was appointed. That representation was unqualified.

Defendant relied on it entirely. She had no basis to monitor for proceedings she
had been told would not occur.

Following that representation, and while Defendant had no counsel and no adequate
notice, the following occurred:

\begin{enumerate}[leftmargin=*, label=\arabic*., itemsep=4pt]

\item The court, apparently on its own, identified that the threshold for
mandatory dismissal under \textit{State v.\ Roberts}, 374 Or.\ 821 (2026), had
been crossed in case 24CR27775. Defendant received no notice of this development.
The first Defendant learned of any dismissal proceeding was upon learning that the
State had filed an objection to it.

\item The State's objection rested on the assertion that Defendant had ``fired''
two lawyers. That assertion is directly refuted by the State's own conduct. At
the hearing on the withdrawal of Defendant's first appointed counsel, the State
initially opposed that withdrawal on the mistaken belief that Defendant had sought
it --- and withdrew its opposition only upon being corrected on the record. An
entity cannot assert that a defendant fired her lawyer while simultaneously
having opposed that lawyer's departure. The State's own on-record conduct at the
Sidman hearing is a judicial admission that Defendant did not cause counsel's
first withdrawal. The withdrawal of Anna Lasher arose from a mandatory conflict
of interest under RPC 1.7, attributed by the Court to counsel's circumstances,
not Defendant's conduct --- stated on the record, April 28, 2025.

\item A hearing was held on the State's objection to the court's sua sponte
Roberts consideration. Defendant was not present. Her remote appearance motion
had been denied. A failure-to-appear notation followed.

\item Show cause proceedings were initiated and an arrest warrant was issued in
case 25CN05790, based on a warrant affidavit whose predicate Defendant contests
as void on its face. Defendant received no adequate notice of those proceedings.

\end{enumerate}

Each of these proceedings was adverse to Defendant's interests. Each was conducted
during a judicially-declared suspension of proceedings. Each proceeded without
the counsel the Court had represented would be obtained before further proceedings
occurred. Defendant's absence from each was not voluntary, was not a waiver of
any right, and was the direct product of the circumstances described in this
Notice.

%% ============================================================
%% SECTION V
%% ============================================================

\section*{V. THE COURT DENIED MOTIONS THE STATE DID NOT OPPOSE --- THEN
LEFT A DEMURRER UNDECIDED FOR FOURTEEN MONTHS}

The April 28, 2025 hearing was convened as an omnibus hearing. Before it was
converted to a withdrawal hearing, prior counsel Anna Lasher's motion to dismiss
and motion to suppress were before the Court. The State had not filed any written
opposition to either motion. Approximately two months after the April 28, 2025
hearing, this Court denied both motions --- without written findings, without a
hearing at which Defendant could be heard, and without any opposition from the
State having ever been filed.

A court that denies motions the opposing party has not contested, without
findings, without hearing, deprives the moving party of a meaningful ruling.
Defendant preserves the denial of those uncontested motions as error.

Separately, Defendant's demurrer --- filed and submitted to this Court on
April 28, 2025 --- has not been decided. It has been pending for more than
fourteen months. ORS 1.050 requires every court to dispose of submitted matters
within three months. That period expired on approximately July 28, 2025. Eleven
months have passed since that date without a ruling, an extension order, or an
explanation.

Under \textit{State v.\ Sanders}, 280 Or.\ 685 (1977), a demurrer must be
resolved before any other pretrial proceedings are conducted. Pretrial proceedings
--- including the adverse hearings described in Section IV --- have been conducted
in this case while the demurrer remained, and remains, undecided. Defendant
preserves all rights arising from that sequence.

Defendant demands a ruling on the demurrer within thirty days of this Notice.

%% ============================================================
%% SECTION VI
%% ============================================================

\section*{VI. THE RECORD DOCUMENTS EXTRA-JUDICIAL COMMUNICATIONS THAT
INITIATED ADVERSE PROCEEDINGS AGAINST DEFENDANT WITHOUT HER KNOWLEDGE}

The record contains documentary evidence --- in the form of a government email
chain produced in discovery in a related proceeding --- establishing the following
sequence:

On June 10, 2025, at 12:20 PM, Christina Stiefer, Pretrial/Verification/Set
Aside/Traffic Supervisor for Washington County Circuit Court, sent an email to
Sheila Clark, P\&P Services Supervisor at Washington County Community Justice.
That email states, in relevant part:

\begin{quote}
``Judge Thompson has forwarded an email to the Pretrial Release Office that you
had sent regarding Ashle Penn. I understand that you have knowledge that Ms. Penn
has been having contact with Mr. Barry Washington who is co-defendant in her
criminal case \# 24CR49043. Her release conditions are to not have any contact
with Mr. Washington. Would you be willing to send our office an email sharing with
us your knowledge of the defendant's contact with her co-defendant? Judge has
requested that we prepare a VRA Show Cause for this violation and I will need a
little more information to support the affidavit.''
\end{quote}

The email chain establishes the following facts, each of which is documented:

\begin{enumerate}[leftmargin=*, label=\arabic*., itemsep=4pt]

\item Sheila Clark of Washington County Community Justice --- a parole and
probation services supervisor whose agency supervised the person identified as
Defendant's co-defendant --- sent an email to the Honorable Brandon Thompson
regarding Defendant.

\item Judge Thompson received that email. He forwarded it to the Washington County
Circuit Court Pretrial Release Office. He directed that office to prepare a VRA
Show Cause proceeding against Defendant.

\item These communications occurred outside the courtroom, outside the record, and
without notice to Defendant or any opportunity for her to be heard.

\item The same day --- June 10, 2025 --- Stiefer received a responsive email from
Molly McBride, a parole and probation officer at Washington County Community
Justice, summarizing contact information gathered through supervised jail calls
placed by the co-defendant. Stiefer filed her Affidavit for Violation of Release
Agreement in case 24CR49043 the same day, June 10, 2025.

\item Circuit Court Judge C. Colburn signed the resulting order directing
Defendant to appear on June 23, 2025.

\item The notice of that June 23, 2025 hearing was directed to 710 NW Autumn
Creek Way, Apt. L304, Beaverton, Oregon --- not Defendant's current address ---
at a time when Defendant's address of record had already been altered as described
in Section III above.

\end{enumerate}

Defendant notes the following for the record. The person identified in the release
agreement as a co-defendant --- Barry Washington --- is a person in whose
supervision proceedings Defendant had been an active and documented advocate,
asserting constitutional challenges to supervision conditions on his behalf
directly to the supervising agency and to oversight bodies. Washington County
Community Justice was aware of Defendant's advocacy. The use of information
gathered through that agency's supervision of the co-defendant, transmitted to
a sitting judge assigned to Defendant's case, outside the record, to initiate
adverse judicial proceedings against Defendant, is a fact the record establishes.
Defendant places it in this record and reserves all rights arising from it.

Or.\ Const.\ Art.\ I, \S~10 guarantees every person a remedy by due course of
law for injury done to them. Canon 2, Rule 2.9 of the Oregon Code of Judicial
Conduct prohibits ex parte communications on pending or impending matters. A
judge who receives correspondence from a parole supervision agency regarding a
defendant in a pending case, forwards that correspondence to his own court's
administrative office, and directs the preparation of adverse proceedings against
that defendant --- all outside the courtroom and without notice to the defendant
--- is not a neutral arbiter of those proceedings or of any proceeding that flowed
from that directive.

Defendant does not ask this Court to rule on this issue here. She asks that the
record reflect what the documents show.

%% ============================================================
%% SECTION VII
%% ============================================================

\section*{VII. THE STIEFER AFFIDAVIT IS PREDICATED ON A VOID STATUTORY
FOUNDATION}

The Affidavit for Violation of Release Agreement filed by Christina Stiefer on
June 10, 2025 invokes ORS 135.290 as the basis for contempt. ORS 135.290
authorizes punishment by contempt of court only for breach of regulations imposed
pursuant to \textbf{ORS 135.260} --- the conditional release statute.

Defendant's release instrument, attached to Stiefer's own affidavit as an
exhibit, is a \textbf{Security Release Agreement} executed under ORS 135.265.
A security release is not a conditional release under ORS 135.260. ORS 135.265
governs release upon posting of security. It does not authorize the imposition of
behavioral conditions as a matter of statutory right. The authority to impose
behavioral conditions through a security release derives from the court's general
authority to set conditions of pretrial release --- not from ORS 135.260.

The contempt remedy under ORS 135.290 requires a conditional release predicate
under ORS 135.260. That predicate does not exist here. The statutory architecture
of the contempt charge collapses on the face of Stiefer's own attached exhibits,
without reference to any fact in dispute.

Defendant preserves this statutory void-on-its-face challenge to the contempt
proceeding in full.

%% ============================================================
%% SECTION VIII
%% ============================================================

\section*{VIII. REQUEST FOR JUDICIAL NOTICE}

Pursuant to ORS 40.060 (Oregon Rule of Evidence 201), Defendant requests that
this Court take judicial notice of the following facts:

\begin{enumerate}[leftmargin=*, label=\arabic*., itemsep=4pt]

\item At timestamp 18:18 of the April 28, 2025 hearing, the Honorable Brandon
Thompson stated on the record: \textit{``it cancels all future hearing dates at
this time.''}

\item At timestamps 19:05 through 19:09 of the same hearing, the court clerk
confirmed Defendant's current address on the record.

\item The OPDC designation entered in Defendant's cases following the April 28,
2025 hearing reflects no assigned attorney and no attorney-client relationship
at any time material to the proceedings described herein.

\item Case 25CN05790 was filed on December 30, 2025.

\item The hearing notice for the June 23, 2025 VRA show cause proceeding in case
24CR49043 was directed to 710 NW Autumn Creek Way, Apt. L304, Beaverton, Oregon
--- not Defendant's current residential address.

\item The arrest warrant in case 25CN05790 was executed at 17548 NW Springville
Road, Apt. 9, Portland, Oregon --- Defendant's current residential address,
which is the same address confirmed on the record at the April 28, 2025 hearing.

\item Defendant's demurrer was submitted to this Court on April 28, 2025 and
remains undecided. The three-month period prescribed by ORS 1.050 expired on
approximately July 28, 2025.

\item Approximately two months after April 28, 2025, this Court denied prior
counsel's motion to dismiss and motion to suppress in case 24CR49043. The State
had filed no written opposition to either motion at the time of denial.

\item Presiding Judge Order No.\ 325 expired January 3, 2025. Presiding Judge
Order No.\ 350 became effective January 31, 2025, with grand jury orientation
beginning that date. No Presiding Judge Order authorized a grand jury to convene
on January 14, 2025, as confirmed in writing by Washington County Trial Court
Administrator Richard Moellmer.

\item Deputy Rambin's supplemental report, dated December 27, 2024, states that
it was possible Defendant was not the suspect in the September 17, 2024 incident.

\item The indictment returned January 14, 2025 does not include the September 17,
2024 incident or its associated valuation, which formed part of the basis for
the original September 18, 2024 arrest and the original charging instrument.

\item Christina Stiefer's June 10, 2025 email to Sheila Clark states: ``Judge
Thompson has forwarded an email to the Pretrial Release Office that you had sent
regarding Ashle Penn\ldots Judge has requested that we prepare a VRA Show Cause
for this violation.''

\item The Affidavit for Violation of Release Agreement filed by Christina Stiefer
on June 10, 2025, in case 24CR49043, attaches a Security Release Agreement
executed under ORS 135.265 as its supporting exhibit. It invokes ORS 135.290 as
the contempt predicate. ORS 135.290 requires a conditional release under ORS
135.260. The attached exhibit is not a conditional release.

\end{enumerate}

%% ============================================================
%% SECTION IX
%% ============================================================

\section*{IX. PRESERVATION OF ALL CLAIMS AND OBJECTIONS}

Defendant expressly preserves, without limitation and without waiver:

\begin{enumerate}[leftmargin=*, label=\arabic*., itemsep=4pt]

\item All challenges to the sufficiency of the January 14, 2025 indictment as
raised in the demurrer submitted April 28, 2025, including: improper joinder
under ORS 135.630(2); failure to state an offense under ORS 135.630(4); legal
bar to the action under ORS 135.630(5); and lack of definiteness under
ORS 135.630(6).

\item All challenges to the validity of the January 14, 2025 grand jury
proceeding, including: the empanelment gap established by the Moellmer
correspondence; the State's failure to disclose Deputy Rambin's December 27,
2024 supplemental report before the grand jury convened; the presentation of
testimony that materially contradicted the sworn September 18, 2024 probable
cause affidavit without correction; and the State's selective omission of the
September 17, 2024 incident after its own officer admitted the identification
was unreliable.

\item All challenges to the probable cause basis for the September 18, 2024
warrantless arrest, including the substantial preliminary showing under
\textit{Franks v.\ Delaware}, 438 U.S.\ 154 (1978), arising from material
misstatements in the probable cause affidavit contradicted by the arresting
officer's own contemporaneous and supplemental reports and by body worn camera
footage of the September 18, 2024 evidence pickup at Target.

\item All suppression claims as to evidence derived from the September 18, 2024
arrest as fruit of a constitutionally defective warrantless arrest.
\textit{State v.\ Hall}, 339 Or.\ 7 (2005); \textit{Wong Sun v.\ United States},
371 U.S.\ 471 (1963).

\item The claim for mandatory dismissal with prejudice under \textit{State v.\
Roberts}, 374 Or.\ 821 (2026), preserved in full. \textit{Roberts} dismissal is
the court's own obligation, triggered by the statutory threshold. The State's
objection --- resting on the claim that Defendant fired two lawyers, directly
refuted by the State's own on-record conduct at the Sidman withdrawal hearing ---
does not defeat that obligation. The clock has run.

\item The facial statutory voidness of the arrest warrant and contempt charge in
case 25CN05790. ORS 135.290 requires a conditional release predicate under
ORS 135.260. The release instrument is a security release under ORS 135.265.
No valid statutory predicate exists. The warrant and contempt charge are void
on their face as a matter of law.

\item All rights arising from the extra-judicial communications described in
Section VI, including rights under Or.\ Const.\ Art.\ I, \S~10, Canon 2,
Rule 2.9 of the Oregon Code of Judicial Conduct, and the due process clause
of the Fourteenth Amendment to the United States Constitution.

\item All rights arising from the defective notice mechanism described in
Section III, including the right to challenge the legal sufficiency of notice
of any proceeding held after December 30, 2025.

\item All rights arising from the denial of the motion to appear remotely and
the subsequent failure-to-appear notation, including any right to challenge
any order or finding entered as a consequence of Defendant's absence from
proceedings she was denied the means to attend.

\item The right to a speedy trial under Or.\ Const.\ Art.\ I, \S~10 and
ORS 135.747, materially impaired by the State's deliberate avoidance of a
preliminary hearing, its manipulation of the grand jury process, its
continuation of prosecution following its own officer's admission of a possible
wrongful arrest, and its conduct of adverse proceedings against an unrepresented
defendant during a judicially-declared suspension of proceedings.

\end{enumerate}

%% ============================================================
%% SECTION X
%% ============================================================

\section*{X. NOTICE OF FORTHCOMING MOTIONS}

This Notice is the first document in a coordinated filing package that Defendant
intends to submit on the papers. The package will include, at minimum:

\begin{enumerate}[leftmargin=*, label=(\arabic*), itemsep=2pt]
\item Brady and Giglio Violation Notice with Demand for Specific Disclosures
--- 24CR49043;
\item Motion for Declaratory Determination of Pure Questions of Law
--- 24CR49043;
\item Renewed and Supplemented Demurrer, Motion to Dismiss with Prejudice,
Motion to Suppress, and Motion for \textit{Franks} Hearing --- 24CR49043;
\item Motion to Compel Production of Grand Jury Transcript --- 24CR49043;
\item Motion for Mandatory Dismissal with Prejudice under \textit{State v.\
Roberts} --- 24CR27775; and
\item Motion to Recall Warrant --- 25CN05790.
\end{enumerate}

Each motion will be accompanied by a proposed order. Each argument will be
grounded in the record this Notice begins to establish.

The record does not present close questions. It presents a prosecution that has
continued past the point the law permits, on a charging instrument returned by
a grand jury with no lawful authority to convene, predicated on an arrest whose
probable cause basis the arresting officer himself placed in doubt in a sworn
supplemental report, sustained by proceedings held in Defendant's absence through
a notice system corrupted by a filing-date address propagation that began the
same day the contempt case was initiated, and initiated in part by extra-judicial
communications between the assigned judge and an adverse agency conducted outside
the record without notice to Defendant.

Defendant is prepared to demonstrate each of these propositions, document by
document, to this Court. She is beginning that demonstration now.

%% ============================================================
%% SIGNATURE
%% ============================================================

\vspace{24pt}
\noindent Respectfully submitted,

\vspace{48pt}
\noindent\rule{2.5in}{0.4pt}\\[4pt]
Ashle\textquotesingle{} Antoinette Penn\\
Defendant, Pro Se\\[6pt]
\noindent Dated: \today\\[6pt]
\textit{Address:} \underline{\hspace{2.5in}}\\[4pt]
\textit{Phone:} \underline{\hspace{2.5in}}\\[4pt]
\textit{Email:} \underline{\hspace{2.5in}}

%% ============================================================
%% CERTIFICATE OF SERVICE
%% ============================================================

\newpage

\begin{center}
\textbf{CERTIFICATE OF SERVICE}
\end{center}

\vspace{12pt}

I, Ashle\textquotesingle{} Antoinette Penn, certify that on \today\ I served a
true and correct copy of the foregoing \textit{Notice of Appearance Pro Se;
Notice of Preservation of Record; Notice of Systematic Obstruction of Access
to Proceedings; and Objection to Characterization of Defendant's Status} upon
the Washington County District Attorney's Office, 150 N.\ First Avenue,
Suite 300, Hillsboro, Oregon 97124, by the following method:

\vspace{12pt}
\noindent $\square$ \ Electronic Filing (eFiling)
\hspace{0.8in}
$\square$ \ U.S.\ Mail, First Class

\vspace{36pt}
\noindent\rule{2.5in}{0.4pt}\\[4pt]
Ashle\textquotesingle{} Antoinette Penn

\end{document}
